\scanpage{11}
Schémas irréductibles, connexes \ldots / noethériens (attention!), loc. noeth /

réduits, intègres ($={}$ irréduct $+$ réduits) normaux \deleted{\unsure{singuliers|réguliers}} (anneaux locaux n.) \deleted{réguliers (id) \illegible{1}}

global $\left\{\begin{aligned}&\text{Morphismes qu. \uline{cpts}}\\&\text{de \uline{type fini}}\end{aligned}\right\}$ a) b) (cas de $V(M),P(M)$)

local $\left\{\begin{aligned}&\text{loc. \uline{de t.f.}}\qquad\text{a) b) c)}\\&\text{loc. \uline{de prés. fi.}}\qquad\text{a) b) }\deleted{\unsure{\text{c}|\text{c)}}}\end{aligned}\right.$

[$f$ \deleted{\unsure{de t.f.}} loc. de t.f. $\Rightarrow\operatorname{diag}_f$ loc. de prés. finie]

\uline{immersions} \qquad a) b) c)

imm. \uline{fermées}, \uline{ouvertes} [\uline{surjectives}] a) b)

sous-schémas ouverts, fermés [transitif, stable par im. réciproque]

\diagram{\begin{tikzcd}[ampersand replacement=\&]
 X\arrow[d]\&X'\arrow[l]\arrow[d,dash]\\
 Y\arrow[r,dash,"\text{cart.}" description]\&Y'
\end{tikzcd}}

$X\overset{f}{\mathrel{-}}Y\overset{g}{\mathrel{-}}Z$

\[
\begin{aligned}
 &\left\{\begin{aligned}
 &\text{a) Si } f\in M,\ f'\in M\\
 &\quad(\Rightarrow\text{ tout isom }\in M)\\
 &\text{b) Si }f,g\in M,\ gf\in M
 \end{aligned}\right.\\
 &\deleted{\unsure{\Leftarrow}}\Longrightarrow(f_1\in M,\ f_2\in M\Longrightarrow f_1\times f_2\in M)\\
 &\Longrightarrow(gf\in M,\ \operatorname{diag}(g)\in M\Longrightarrow f\in M)\quad *
\end{aligned}
\]

c) Si $gf\in M$, $f\in M$ [en présence de a) et b), équivaut à : toute section est $\in M$, \{ou : diagonales $\in M$\} ; ou : noyaux $\in\underline{M}$ si c'est local en bas \inserted{et vrai loclt sur \unsure{but}} \} on est ramené au cas affine]

\margin{Ex : vérifier sur diagonal, sections, noyaux}

morphismes \uline{séparés} [a) b) c)] : diagonale fermée $\Longleftrightarrow$ noyaux de couples dedans fermés $\Longleftrightarrow$ principe d'identité pour morphismes des schémas \uline{réduits} dedans.

\uline{Ex} (mono -- donc applic. diag, noyaux, diag. morphismes affines

\deleted{\unsure{morphismes quasi-séparés}}

$P(V)\longrightarrow S$

Schémas séparés (ils sont séparés sur n'importe quoi). $\Rightarrow$ \uline{Intersection} \deleted{\unsure{affines}} de deux ouverts affines affine sur $S$ i.e. $X\unsure{,}Y$ sur $S$-affine $\Rightarrow X\times_S Y$ affine

% slanted
\margin{on \uline{peut} \uline{relativiser} $gf$ affine, $g$ séparé $\Rightarrow f$ affine}

morphismes \uline{quasi-séparés} (a) b) c)) (\deleted{\unsure{not}} schémas quasi-séparés)

diagonale quasi-cpte $\Longleftrightarrow$ noyaux de couples dedans quasi-compacts $\Longleftrightarrow$ sections après chgt de base qu. cpt

\deleted{\unsure{et la base}} [Si base qu. sép : $\Longleftrightarrow$ \deleted{\unsure{la diagonale qu. cpt}} l'intersection de 2 ouverts qu. cpts est \unsure{enc. q|qu. cpte}

\deleted{\unsure{morphismes} \unsure{séparés} (\unsure{cf.}) \unsure{et qu. ctp et quasi-séparés})}

$\Longleftrightarrow$ \deleted{tout} schém. \unsure{sur} $X$ qu. cpt. est qu. compact sur $X$.]

% slanted
\margin{$gf$ qu. cpt (de t.f.) $g$ qu. sép $\Rightarrow$ $f$ qu. cpt (qu. sép)}

morphismes \uline{cohérents} $={}$ qu. cpt qu. sép \quad a) b) [$gf$ coh et $g$ qu. sép $\Rightarrow f$ coh]

Si $f:X\mathrel{-}Y$ qu. coh. et $F$ qu. coh. sur $X$, alors $f_*(F)$ (et les $R^i f_*(F)$) quasi-cohérent.

\margin{schémas cohérents}

\uline{NB} Tous les schémas envisagés par les old-timers étaient de t.f. et \uline{séparés} sur des corps. La plupart des schémas qu'on envisage \unsure{encore aujourd'hui|encore auj.|couramment} sont séparés pratiquement \uline{tous} ceux qu'on rencontre sont quasi-séparés.

Beaucoup sont qu. cpts -- mais il y en a d'importants qui ne le sont pas.

[morphismes de \uline{prés. finie} ($={}$ loc. de prés. finie et cohérents) \quad a) b) \quad [$gf$ de prés. finie et $g$ de t.f et qu. sép.

\fbox{Toutes les notions sont \uline{locales en bas}}

$\longrightarrow$ Yoga des passages à la limite.

\uline{Généralités} \quad Local en bas $\left\{\begin{aligned}&\text{stable par image inverse}\\&\text{stable par composition}\end{aligned}\right.$ $\substack{\Longrightarrow\\\Longrightarrow}$

$g:Y\mathrel{-}S$

\[
\begin{aligned}
\operatorname{diag}g\in M &\Longleftrightarrow \forall S'\mathrel{-}S,\ \text{tte section de }Y'\text{ sur }S'\text{ est }\in M\\
 &\Longleftrightarrow \forall (u,v):X\rightrightarrows Y\text{ sur }S,\quad\ker(u,v)\lhook\joinrel\longrightarrow X\text{ est }\in M\\
 &\Longleftrightarrow \forall f:X\mathrel{-}Y,\quad\Gamma_f:X\mathrel{-}X\times_S Y\text{ est }\in M\\
 &\Longleftrightarrow gf\in M\Longrightarrow f\in M
\end{aligned}
\]

\diagram{\begin{tikzcd}[ampersand replacement=\&,column sep=small]
 *\quad X\arrow[d,"f"']\&{}\\
 Y\arrow[d,"g"']\&X\times_S Y\arrow[l,"\mathrm{pr}_2"']\arrow[d,"\mathrm{pr}_1"]\\
 S\&X\arrow[l,"gf"]\arrow[u,bend right,"\Gamma_f"']
\end{tikzcd}}

\diagram{\begin{tikzcd}[ampersand replacement=\&]
 X\arrow[r,"\Gamma_f"]\arrow[d,"f"']\&X\times_S Y\arrow[d,"f\times\mathrm{id}_Y"]\\
 Y\arrow[r,hook,"\operatorname{diag}_Y"']\&Y\times_S Y
\end{tikzcd}}

\uline{plan} \quad a) Traversée du désert

b) $\left\{\begin{aligned}&\text{Passage au quotient}\\&\text{descente fid. plat}\end{aligned}\right\}$ et langage faisceautique

c) Notions ``géométriques'' sur un corps

d) \deleted{\illegible{1}} Ensembles constructibles, passages à la limite

% archive sheet number: 10
